Recent studies suggest that large language model (LLM) hallucinations may influence molecular property prediction, yet the nature of this effect remains poorly characterized. We introduce a controlled semantic perturbation framework to evaluate how structured modifications --- which we characterize as synthetic semantic priors --- influence zero-shot prediction. We evaluate nine conditions on two MoleculeNet benchmarks (BBBP, $N=204$; BACE, $N=152$) using multi-seed evaluation (Llama-3.1-8B) and cross-model validation (Llama-3.3-70B). On BBBP, prompt-constrained structural augmentation (C4a) yields the largest observed ROC-AUC increase (0.626 vs. 0.533 baseline), although shifts were not statistically significant after correction ($p > 0.05$). Post-hoc verification suggests that the structural topic focus of the perturbation, rather than its factual accuracy, is the primary driver of performance shifts, as C4a outputs predominantly fall into contextual descriptions despite their structural focus. On BACE, we find robust and statistically significant performance degradation across nearly all hallucination conditions, ditandai oleh systematic ``distributional compression'' of predicted probabilities toward a non-discriminative low-confidence state. A random-permutation control confirms that observed effects are driven by semantic alignment rather than scientific register alone, suggesting a functional link between perturbation content and model response. These findings indicate that LLM-based molecular predictions exhibit differential sensitivity to structured semantic perturbations, highlighting the importance of controlled evaluation in scientific reasoning pipelines.
